\section{Known quadratic trace identities}\label{app:quadratic}
Let $\chi$ be the quadratic character of $\F_p$, extended by zero, and put $T_j(p)=\sum_{\lambda\ne0,1}\chi(\lambda)a_p(\lambda)^j$. The following is a known-trace-formula corollary.
\begin{theorem}\label{thm:quadratic}
For every prime $p\ge5$,
\[
 T_4(p)=-c(p)-b(p)-3p\,a(p)-2p^2-3p-1,
\]
where $a=8.4.a.a=\eta(2z)^4\eta(4z)^4$,
$b=4.6.a.a=\eta(2z)^{12}$, and
$c=8.6.a.a=q+20q^3-74q^5-24q^7+\cdots$.
\end{theorem}
\begin{proof}
Use the Ahlgren and Ahlgren--Ono trace formulas as stated in
\cite[Theorem 2.3]{Grove}:
\begin{align*}
 -\Tr_{k+2}(\Gamma_0(4),p)&=3+\sum_{\lambda\ne0,1}G_{k+2}(a(\lambda),p),\\
 -\Tr_{k+2}(\Gamma_0(8),p)&=4+\sum_{\gamma\ne0,1,-1}G_{k+2}(a(\gamma^2),p).
\end{align*}
Here $G_4(s,p)=s^2-p$ and $G_6(s,p)=s^4-3ps^2+p^2$.
Put $U_k=\sum_{\lambda\ne0,1}a(\lambda)^k$ and
$Q_k=\sum_{\gamma\ne0,1,-1}a(\gamma^2)^k$. Then
$Q_k-U_k=T_k$, and $Q_0-U_0=-1$.
The relevant spaces have bases
\[
 S_4(\Gamma_0(4))=0,\quad S_4(\Gamma_0(8))=\langle a\rangle,
\quad S_6(\Gamma_0(4))=\langle b\rangle,
\]
\[
 S_6(\Gamma_0(8))=\langle b(z),b(2z),c(z)\rangle.
\]
Subtracting the weight-four trace formulas gives
$-a(p)=1+T_2(p)+p$, so $T_2=-a(p)-p-1$.
Subtracting the weight-six formulas gives
$-(b(p)+c(p))=1+T_4-3pT_2-p^2$.
Substitution proves the formula.
\end{proof}

The identity is a consequence of established trace formulas. No new
modularity or priority assertion is attached to this appendix.

\section{The corrected hypergeometric normalization}\label{sec:hyp}
Assume $p\ge5$, $p\equiv1\pmod3$. Let $\phi$ and $\rho$ be the quadratic and a primitive cubic character, and let $g$ and $J$ denote Gauss and Jacobi sums for a fixed nontrivial additive character. We use the normalized Gauss-sum definition
\cite[Eq. (1.1)]{Grove}, with all multiplicative characters, including
the trivial character, extended by zero at zero. Put
\[
 H(\lambda)=\frac1{1-p}\sum_\chi c(\chi)\chi(\lambda),
 \qquad c(\chi)=\frac{g(\phi\chi)^2g(\bar\chi)^2}{g(\phi)^2}.
\]
Only nonzero arguments are used here; the separately defined value at zero is immaterial. For nonsingular $\lambda$, $a(\lambda)=\phi(-1)H(\lambda)$.
Fix a generator $\omega$ of the multiplicative character group such that $\rho=\omega^{(p-1)/3}$ and $\phi=\omega^{(p-1)/2}$. The same $\omega$ is used to interpret the rational parameters of $H_p(A;1)$. Let
\[
 A=\left\{\tfrac12,\tfrac12,\tfrac13,\tfrac13;
               1,1,\tfrac56,\tfrac56\right\},\qquad
 J=J(\rho,\phi\bar\rho).
\]
\begin{lemma}\label{lem:norm}
With this normalization,
\begin{equation}\label{eq:correctnorm}
 M_{3,2}(p)+1
 =\frac1{p-1}\sum_\chi c(\chi)c(\bar\rho\bar\chi)
 =-J^2 H_p(A;1).
\end{equation}
\end{lemma}
\begin{proof}
Since $\phi$ is nontrivial,
$c(\chi)=J(\phi\chi,\bar\chi)^2$ by the Gauss--Jacobi identity.
Expanding the Jacobi sums and summing over $\chi$ imposes
$xy=(1-x)(1-y)$, equivalently $x+y=1$. Therefore
\[
 \sum_\chi c(\chi)=(p-1)\sum_x\phi(x(1-x))
 =-(p-1)\phi(-1),
\]
so $H(1)=\phi(-1)$. Character orthogonality applied to $H(\lambda)^2$
then gives the first equality of \eqref{eq:correctnorm}.

The numerator of each summand for $H_p(A;1)$ is
\[
 g(\phi\chi)^2g(\rho\chi)^2g(\bar\chi)^2
                  g(\phi\bar\rho\bar\chi)^2.
\]
Its normalizing denominator is
$g(\phi)^2g(\rho)^2g(\varepsilon)^2g(\phi\bar\rho)^2$,
where $g(\varepsilon)^2=1$. The convolution on the middle of
\eqref{eq:correctnorm} has denominator $g(\phi)^4$ and the opposite
overall sign, since $1-p=-(p-1)$. Their ratio is
\[
 -\left(\frac{g(\rho)g(\phi\bar\rho)}{g(\phi)}\right)^2=-J^2.
\]
This proves the normalization factor.
\end{proof}

For instance, choose $\omega(3)=e^{2\pi i/6}$ at $p=7$ and
$\rho=\omega^2$. Then $J=2+\sqrt{-3}$, $J^2=1+4\sqrt{-3}$,
and $M_{3,2}(7)+1=-15$. Hence
$H_7(A;1)=15/(1+4\sqrt{-3})$, which is not $-15$.
The v1 phrase identifying the convolution directly with the normalized
$A$-value must therefore include the factor in \eqref{eq:correctnorm}.
The underlying moment calculation is unchanged.

Under the carried-certificate comparison of Section~\ref{sec:cubic}, the corrected modular relation is
\[
 J^2 H_p(A;1)=a_f(p)+p.
\]
For $B=\{\frac12,\frac12,\frac13,\frac23;1,1,1,1\}$, the
published identity $H_p(B;1)=a_f(p)+\leg{3}{p}p$
\cite[Theorem~2 and Table~1]{LTYZ} gives
\[
 J^2 H_p(A;1)=H_p(B;1)+\left(1-\leg{3}{p}\right)p.
\]
The datum $A$ also occurs in the EHMM literature at a different
argument \cite{EHMM2}. A priority claim requires comparison of both
the argument and the Jacobi twists. The algebra above supplies the normalization proof; numerical scripts are retained separately as historical checks.

